\documentclass[conference]{IEEEtran}

% ---------------- Packages ----------------
\usepackage{graphicx}
\usepackage{amsmath}
\usepackage{booktabs}
\usepackage{hyperref}
\usepackage{float}

% ---------------- Title ----------------
\title{A Modular Machine Learning Portfolio Dashboard Using Streamlit for Interactive Prediction and Analysis}

\author{
\IEEEauthorblockN{Karim Guida}
\IEEEauthorblockA{
Ottawa, Canada \\
Email: guida.karim@gmail.com
}
}

\begin{document}

\maketitle

% ---------------- Abstract ----------------
\begin{abstract}
This paper presents the design and implementation of a modular machine learning portfolio application developed using Streamlit. The system integrates four distinct machine learning tasks—classification, regression, clustering, and neural network-based prediction—into a unified, interactive dashboard. Each module follows a structured pipeline including data preprocessing, model training, evaluation, and deployment. The application emphasizes usability, modularity, and real-time inference, providing a practical demonstration of applied machine learning techniques across multiple domains.
\end{abstract}

\begin{IEEEkeywords}
Machine Learning, Streamlit, Classification, Regression, Clustering, Neural Networks, Data Science, Interactive Dashboard
\end{IEEEkeywords}

% ---------------- Introduction ----------------
\section{Introduction}

Machine learning applications are increasingly deployed in interactive environments to facilitate decision-making and exploratory analysis. Traditional workflows often separate data analysis, model training, and deployment, resulting in fragmented solutions.

This project aims to bridge this gap by developing a unified, modular dashboard that integrates multiple machine learning models into a single interactive application. The system is designed to demonstrate real-world machine learning use cases while maintaining a clean architecture and user-friendly interface.

% ---------------- System Architecture ----------------
\section{System Architecture}

The application follows a modular architecture where each machine learning project is encapsulated within its own module. The system is structured as follows:

\begin{itemize}
    \item \textbf{Data Layer:} Raw datasets stored in CSV format
    \item \textbf{Model Layer:} Trained models serialized using pickle
    \item \textbf{Logic Layer:} Prediction and preprocessing functions
    \item \textbf{Presentation Layer:} Streamlit-based user interface
\end{itemize}

This separation ensures scalability, maintainability, and ease of extension.

% ---------------- Methodology ----------------
\section{Methodology}

Each module follows a consistent pipeline:

\subsection{Data Preprocessing}
Data is cleaned, transformed, and encoded as required. Numerical features are scaled where necessary, and categorical variables are converted using encoding techniques.

\subsection{Model Training}
Different models are used depending on the problem:

\begin{itemize}
    \item Classification: Random Forest
    \item Regression: Random Forest Regressor
    \item Clustering: K-Means
    \item Neural Network: Multi-Layer Perceptron (MLPRegressor)
\end{itemize}

\subsection{Model Evaluation}
Models are evaluated using appropriate metrics such as accuracy for classification and RMSE for regression.

\subsection{Deployment}
Trained models are saved and integrated into the Streamlit application for real-time inference.

% ---------------- Case Studies ----------------
\section{Case Studies}

\subsection{Loan Eligibility Prediction}
A classification model predicts whether a loan application will be approved based on applicant features such as income, credit history, and education.

\subsection{Real Estate Price Prediction}
A regression model estimates housing prices using structural and financial attributes of properties.

\subsection{Customer Segmentation}
K-Means clustering is applied to group customers based on income and spending behavior. The model identifies distinct customer segments for targeted analysis.

\subsection{Admission Prediction}
A neural network model predicts graduate admission probability using academic indicators such as GRE scores, TOEFL scores, and CGPA.

% ---------------- Results ----------------
\section{Results}

The application successfully integrates multiple models into a single interface, enabling:

\begin{itemize}
    \item Real-time prediction
    \item Interactive data exploration
    \item Visualization of clustering results
    \item Consistent user experience across modules
\end{itemize}

The modular design allows each component to function independently while maintaining overall system cohesion.

% ---------------- Discussion ----------------
\section{Discussion}

The project demonstrates the effectiveness of combining machine learning models with interactive dashboards. The use of Streamlit simplifies deployment while maintaining flexibility.

Challenges encountered include:
\begin{itemize}
    \item Ensuring consistent UI styling across components
    \item Handling different model types within a unified interface
    \item Managing feature consistency between training and inference
\end{itemize}

% ---------------- Conclusion ----------------
\section{Conclusion}

This work presents a modular machine learning portfolio that integrates multiple predictive models into a cohesive and interactive application. The system highlights the importance of clean architecture, usability, and real-time inference in modern data science applications.

Future work may include cloud deployment, API integration, and enhanced visualization capabilities.

% ---------------- References ----------------
\begin{thebibliography}{00}

\bibitem{b1} F. Pedregosa et al., ``Scikit-learn: Machine Learning in Python,'' Journal of Machine Learning Research, 2011.

\bibitem{b2} Streamlit Inc., ``Streamlit Documentation,'' \url{https://streamlit.io}

\bibitem{b3} T. Hastie, R. Tibshirani, J. Friedman, ``The Elements of Statistical Learning,'' Springer, 2009.

\end{thebibliography}

\end{document}